\documentclass[a4paper,11pt]{article}

\usepackage[margin=2cm]{geometry}
\usepackage{array}
\usepackage{tabularx}
\usepackage{multirow}
\usepackage{graphicx}
\usepackage{rotating}
\usepackage{setspace}

\renewcommand{\arraystretch}{1.3}

\newcolumntype{Y}{>{\raggedright\arraybackslash}X}

\begin{document}
    
    \section{Assessment Criteria ``Writing''}
        
        In the subtest ``Writing'', the writing performance of the participants is evaluated.
        A distinction is made between content-related and language-related criteria.
        While the first evaluates how well the task has been completed, the language-related
        criteria are based on the Common European Framework of Reference for Languages.
        
        The following criteria are considered:
        
        \subsection{Content Appropriateness}
            
            \begin{enumerate}
                \item Task completion
            \end{enumerate}
        
        \subsection{Language Appropriateness}
            
            \begin{enumerate}
                \item Communicative design
                \item Correctness
                \item Vocabulary
            \end{enumerate}
        
        \subsection{Content Appropriateness}
            
            \begin{center}
                \small
                \begin{tabularx}{\textwidth}{|c|Y|Y|Y|Y|Y|}
    \hline
    & \textbf{5 Points} & \textbf{4 Points} & \textbf{3 Points} & \textbf{2 Points} & \textbf{1 Point} \\
    \hline
    
    \rotatebox{90}{Task Completion}
    
    &
    All four guiding points are addressed precisely.
    
    &
    All four guiding points are addressed, reader cooperation required.
    Or: three guiding points appropriately addressed.
    
    &
    Three guiding points addressed, reader cooperation required.
    Or: two guiding points appropriately addressed.
    
    &
    Two guiding points addressed, reader cooperation required.
    Or: one guiding point appropriately addressed.
    
    &
    Only one guiding point addressed, reader cooperation required.
    
    \\
    \hline
                \end{tabularx}
            \end{center}
            
            \newpage
    
    
    \section{Language Appropriateness}
        
        \begin{center}
            \footnotesize
            \begin{tabularx}{\textwidth}{|c|Y|Y|Y|}
    \hline
    
    & \textbf{B1} & \textbf{A2} & \textbf{A1} \\
    
    \hline
    
    \rotatebox{90}{Communicative Design}
    
    &
    Can realize a wide range of language functions and respond to them using common expressions.
    
    Can link short elements into a simple connected statement.
    
    &
    Can use basic language functions such as exchanging information, making requests,
    expressing opinions and attitudes.
    
    Can connect simple sentences using common connectors such as ``and'', ``but'', ``because''.
    
    &
    Can establish basic social contact using simple greetings and expressions.
    
    Can connect words or phrases with simple connectors such as ``and'' or ``then''.
    
    \\
    \hline
    
    \rotatebox{90}{Correctness}
    
    &
    Generally good control of grammar despite noticeable influence of the native language.
    Errors occur but meaning remains clear.
    
    Spelling and punctuation are mostly understandable.
    
    &
    Can use some structures correctly but still makes systematic basic errors such as tense confusion
    or missing subject–verb agreement.
    
    Can write words from spoken vocabulary approximately correctly.
    
    &
    Shows limited control of a few simple grammatical structures and memorized sentence patterns.
    
    Can copy familiar words and short phrases.
    
    \\
    \hline
    
    \rotatebox{90}{Vocabulary}
    
    &
    Has a sufficient vocabulary to describe most everyday situations using some paraphrasing.
    
    Shows good control of basic vocabulary but still makes elementary mistakes when expressing
    more complex ideas.
    
    &
    Has enough vocabulary for routine everyday situations and familiar topics.
    
    Uses a limited vocabulary related to concrete everyday needs.
    
    &
    Uses isolated words and expressions related to specific concrete situations.
    
    Uses individual words and short sentences for basic needs.
    
    \\
    \hline
            \end{tabularx}
        \end{center}

\end{document}